%%%%%%%%%%%%%%%%%%%%%%%%%%%%%%%%%%%%%%%%%%%%%%%%%%%%%%%%%%%%%%%%%%%%%%
% Lớp 8 | Chương III. Tứ giác | Hình chữ nhật
% Dạng 8H3?13-2: Dấu hiệu nhận biết hình chữ nhật
%%%%%%%%%%%%%%%%%%%%%%%%%%%%%%%%%%%%%%%%%%%%%%%%%%%%%%%%%%%%%%%%%%%%%%


% nguồn: De_AI_Toan8_cuoi_hoc_ki_i_1 (nhập 21/11/2026)
\begin{bt}%[8H3H13-2]%[ma:4f930a69]
	Cho tam giác $ABC$ vuông tại $A$. Lấy điểm $M$ thuộc cạnh $BC$. Kẻ $MD \perp AB$ tại $D$ và $ME \perp AC$ tại $E$.
a) Chứng minh tứ giác $ADME$ là hình chữ nhật.
b) Chứng minh $AM = DE$.
	\loigiai{
		a) Xét tứ giác $ADME$ có: \\
$\widehat{DAE} = 90^\circ$ (tam giác $ABC$ vuông tại $A$); \\
$\widehat{ADM} = 90^\circ$ ($MD \perp AB$); \\
$\widehat{AEM} = 90^\circ$ ($ME \perp AC$). \\
Tứ giác $ADME$ có ba góc vuông nên $ADME$ là hình chữ nhật. \\
b) Vì $ADME$ là hình chữ nhật nên hai đường chéo của nó bằng nhau, suy ra $AM = DE$.
	}
\end{bt}

% nguồn: De_AI_Toan8_cuoi_hoc_ki_i_1 (nhập 21/11/2026)
\begin{bt}%[8H3H13-2]%[ma:f7c6fe60]
	Cho tam giác $ABC$ cân tại $A$ có đường cao $AH$ ($H \in BC$). Gọi $I$ là trung điểm của $AC$, lấy điểm $K$ đối xứng với $H$ qua $I$.
a) Chứng minh tứ giác $AHCK$ là hình chữ nhật.
b) Chứng minh $AK \parallel BC$ và $AK = BH$.
	\loigiai{
		a) Vì $I$ là trung điểm của $AC$ và $K$ đối xứng với $H$ qua $I$ nên $I$ cũng là trung điểm của $HK$. \\
Tứ giác $AHCK$ có hai đường chéo $AC$ và $HK$ cắt nhau tại trung điểm $I$ của mỗi đường nên $AHCK$ là hình bình hành. \\
Lại có $\widehat{AHC} = 90^\circ$ ($AH \perp BC$), nên hình bình hành $AHCK$ là hình chữ nhật. \\
b) Vì $AHCK$ là hình chữ nhật nên $AK \parallel HC$ (tức $AK \parallel BC$) và $AK = HC$. \\
Vì tam giác $ABC$ cân tại $A$ nên đường cao $AH$ cũng là đường trung tuyến, suy ra $BH = HC$. Do đó $AK = BH$.
	}
\end{bt}

% nguồn: De_AI_Toan8_cuoi_hoc_ki_i_1 (nhập 21/11/2026)
\begin{bt}%[8H3H13-2]%[ma:ed4a52eb]
	Cho tam giác $ABC$ vuông tại $A$, gọi $M$ là trung điểm của cạnh $BC$. Lấy điểm $D$ đối xứng với $A$ qua $M$.
a) Chứng minh tứ giác $ABDC$ là hình chữ nhật.
b) Biết $AB = 5\text{ cm}$, $AC = 12\text{ cm}$, tính độ dài đoạn thẳng $AD$.
	\loigiai{
		a) Tứ giác $ABDC$ có hai đường chéo $AD$ và $BC$ cắt nhau tại trung điểm $M$ của mỗi đường nên $ABDC$ là hình bình hành. \\
Lại có $\widehat{BAC} = 90^\circ$ (tam giác $ABC$ vuông tại $A$), nên hình bình hành $ABDC$ là hình chữ nhật. \\
b) Áp dụng định lí Pythagore trong tam giác $ABC$ vuông tại $A$: $BC^2 = AB^2 + AC^2 = 5^2 + 12^2 = 25 + 144 = 169 \Rightarrow BC = 13\text{ cm}$. \\
Vì $ABDC$ là hình chữ nhật nên hai đường chéo bằng nhau: $AD = BC = 13\text{ cm}$.
	}
\end{bt}

% nguồn: De_AI_Toan8_cuoi_hoc_ki_i_1 (nhập 21/11/2026)
\begin{bt}%[8H3H13-2]%[ma:ad39a3f1]
	Cho hình bình hành $ABCD$ (không phải hình thoi). Các tia phân giác của bốn góc $\widehat{A}$, $\widehat{B}$, $\widehat{C}$, $\widehat{D}$ cắt nhau tạo thành một tứ giác $MNPQ$ (với $M$ là giao điểm phân giác góc $A$ và $B$, $N$ là giao điểm phân giác góc $B$ và $C$, $P$ là giao điểm phân giác góc $C$ và $D$, $Q$ là giao điểm phân giác góc $D$ và $A$). Chứng minh tứ giác $MNPQ$ là hình chữ nhật.
	\loigiai{
		Trong hình bình hành $ABCD$, ta có $\widehat{A} + \widehat{B} = 180^\circ$ (hai góc trong cùng phía). \\
Vì $AM$, $BM$ là tia phân giác của $\widehat{A}$, $\widehat{B}$ nên $\widehat{MAB} + \widehat{MBA} = \dfrac{\widehat{A} + \widehat{B}}{2} = \dfrac{180^\circ}{2} = 90^\circ$. \\
Xét tam giác $AMB$, ta có $\widehat{AMB} = 180^\circ - 90^\circ = 90^\circ$, suy ra $\widehat{QMN} = 90^\circ$ (đối đỉnh hoặc kề bù tùy vị trí, đều bằng $90^\circ$). \\
Chứng minh hoàn toàn tương tự với các cặp góc kề $(\widehat{B}, \widehat{C})$ và $(\widehat{C}, \widehat{D})$, ta được $\widehat{MNP} = 90^\circ$ và $\widehat{NPQ} = 90^\circ$. \\
Tứ giác $MNPQ$ có ba góc vuông nên là hình chữ nhật.
	}
\end{bt}
